\subsection{Questão 28}

\subsubsection{Enunciado}

A Fig.~\ref{fig:35-42} mostra duas fontes luminosas isotrópicas, \(S_1\) e \(S_2\), que emitem em fase com um comprimento de onda de \(400\,\text{nm}\) e mesma amplitude.

Um detector \(P\) é colocado sobre o eixo \(x\), que passa pela fonte \(S_1\).

A diferença de fase \(\phi\) entre os raios provenientes das duas fontes é medida entre \(x = 0\) e \(x = +\infty\); os resultados entre \(0\) e
\[
x_s = 10 \times 10^{-7}\,\text{m}
\]
aparecem na Fig.~35-42.

Qual é o maior valor de \(x\) para o qual os raios chegam ao detector \(P\) com fases opostas?

\begin{figure}[h]
    \centering
    \imagemquestao{figures/fig35-42.jpeg}
    \label{fig:35-42}
\end{figure}

\vspace{0.5cm}

\subsubsection{Resolução}


\vspace{1cm}
